\chapter{Mathematical Logic}\label{logic}
\chapterstatus{planned}

The formal language in which the rest of the project is written: what a statement, a proof and a theory are, and what can and cannot be expected of them.

\noindent\textbf{Prerequisites.} None.

\noindent\textbf{Planned contents.}
\begin{itemize}
\item Propositional logic: formulas, truth tables, tautologies, normal forms.
\item First-order languages: terms, formulas, free and bound variables, substitution.
\item Structures and satisfaction; models of a theory; semantic consequence.
\item Formal proofs and the soundness theorem.
\item G\"odel's completeness theorem and the compactness theorem.
\item The L\"owenheim--Skolem theorems and non-standard models.
\item Statements of G\"odel's incompleteness theorems and what they mean for foundations.
\end{itemize}
